\chapter*{Abstract}
\addcontentsline{toc}{chapter}{Abstract}

When gravitational-wave (GW) detectors flag a candidate compact-binary merger, full parameter estimation can take hours to days. In that window, the only public information is the stream of GCN Circulars, free-text bulletins issued by the LIGO/Virgo/KAGRA collaboration and by follow-up observers. This project investigates whether that text alone can support a fast, approximate estimate of chirp mass -- the single parameter most informative about a merger's astrophysical nature -- before full parameter estimation completes.

A reproducible pipeline was built to convert the raw circular archive (45{,}268 circulars) into a structured, per-event dataset with field-level provenance. Applied to the O4a observing run (181 events), the resulting data-availability audit shows that two of the three inputs a standard SNR-scaling chirp-mass formula requires -- network signal-to-noise ratio and orbital inclination -- are structurally absent from circular text, not merely difficult to parse. A physics-based estimator built on that formula was tested against 49 events where circulars self-report a reference chirp mass, first with the formula's literature-quoted calibration constant and then with that constant properly fit via leave-one-out cross-validation; calibration improves the median error substantially (115\% to 45\%) but leaves a statistically significant residual bias correlated with distance ($r = 0.55$), showing the shortfall is structural rather than a poorly chosen constant. A statistical alternative -- a log-linear fit of chirp mass on distance alone, requiring no SNR assumption -- performs better (33\% median error, correct mass-bin 51\% of the time); a trivial source-class-midpoint baseline is competitive with it on several metrics, though this is better explained by the validation sample being 98\% one source class than by any genuine distance-class relationship.

This structural diagnosis was then tested directly rather than left as an inference: cross-matching circular event IDs against the public GWTC catalog (via GWOSC) gives a 214-event set with real, measured network SNR -- the one input circulars cannot supply. Holding every other input and procedure fixed and substituting real SNR for the assumed value cuts the physics formula's median error from 46.6\% to 31.5\% and nearly doubles the fraction of events within $\pm25\%$ (27\% to 42\%). This directly confirms, rather than merely argues, that SNR was the dominant error source: the formula is not unsupported in the abstract, only as a circular-derived estimator specifically.

The thesis further evaluates general-purpose large language models on GW-domain information-extraction tasks, finding inconsistent reliability from unaided model memory, and uses that finding to motivate the design (though not yet a full implementation) of a retrieval-grounded query tool, GW Pathfinder, that answers only from retrieved structured records rather than free generation.

\textbf{Keywords:} gravitational waves, chirp mass, GCN circulars, parameter estimation, information extraction, large language models
